%==============================================================
%  Figura 6.11 - Modello a V semplificato per lo sviluppo del
%  software correlato alla sicurezza
%  Adattamento in italiano della Figura 6.11 dell'IFA Report 2/2017
%
%  Le coordinate sono state ricavate dalla grafica vettoriale
%  dell'originale (pagina 65 di rep0217e.pdf) e sono espresse nei
%  punti PostScript di quel disegno; l'unita' della tikzpicture
%  (x=y=0.042cm) riporta il tutto alle proporzioni originali.
%
%  Sistema di riferimento: l'origine e' l'angolo in basso a sinistra
%  della cornice esterna; l'asse y e' rivolto verso l'alto.
%
%  Il ramo sinistro della "V" (progettazione) scende da "Specifica"
%  a "Codifica"; quello destro (test) risale fino a "Validazione".
%  Le frecce continue sono i risultati di una fase, quelle punteggiate
%  le verifiche che risalgono alla fase corrispondente.
%
%  Il corpo e' \footnotesize e non \small come nella Figura 6.10:
%  le diciture italiane sono piu' lunghe delle inglesi e con \small
%  uscirebbero dai riquadri, che qui hanno le dimensioni dell'originale.
%==============================================================
\begin{tikzpicture}[
  x=0.042cm, y=0.042cm,
  font=\footnotesize,
  cornice/.style={draw=black, line width=0.9pt},
  intestazione/.style={draw=black, fill=black!11, line width=0.9pt},
  blocco/.style={draw=black, fill=black!9, line width=0.9pt,
                 rounded corners=4pt, align=center, inner sep=2pt},
  risultatoIO/.style={draw=black, line width=2.12pt,
                      -{Triangle[length=9.6pt,width=9.6pt]}},
  risultato/.style={draw=black, line width=1.59pt,
                    -{Triangle[length=8.2pt,width=8.2pt]}},
  verifica/.style={draw=black, line width=1.59pt, line cap=butt,
                   dash pattern=on 1.6pt off 1.3pt,
                   -{Triangle[length=8.2pt,width=8.2pt]}},
  diagonale/.style={draw=black, line width=0.9pt,
                    -{Triangle[length=7pt,width=6pt]}}
]

\setstretch{1}

% --- riquadro di intestazione -------------------------------------------
\draw[intestazione] (12.09,249.94) rectangle (407.78,283.91);
\node[align=center, font=\footnotesize\bfseries] at (209.94,266.9)
  {Modello di sviluppo: modello a V semplificato\\
   Obiettivo: software leggibile, comprensibile, testabile e manutenibile};

% --- riquadri delle fasi -------------------------------------------------
% ramo di progettazione (sinistra)
\node[blocco, minimum width=3.368cm, minimum height=2.404cm] (spec)
  at (121.61,194.81) {Specifica del\\software correlato\\alla sicurezza};
\node[blocco, minimum width=2.363cm, minimum height=1.380cm] (progsis)
  at (138.90,131.30) {Progettazione\\del sistema};
\node[blocco, minimum width=2.364cm, minimum height=1.381cm] (progmod)
  at (168.14,79.98) {Progettazione\\dei moduli};
\node[blocco, minimum width=3.148cm, minimum height=1.381cm] (codifica)
  at (209.46,28.58) {Codifica};
% ramo di test (destra)
\node[blocco, minimum width=3.595cm, minimum height=2.407cm] (valid)
  at (302.81,194.78) {Validazione};
\node[blocco, minimum width=2.695cm, minimum height=1.380cm] (testint)
  at (275.26,131.30) {Test di\\integrazione};
\node[blocco, minimum width=2.363cm, minimum height=1.381cm] (testmod)
  at (250.71,79.98) {Test dei\\moduli};

% --- ingresso e uscita ----------------------------------------------------
\draw[risultatoIO] (23.32,198.91) -- (81.51,198.91);
\node[anchor=north west, align=left] at (19.97,235.38)
  {Specifica delle\\funzioni di\\sicurezza};

\draw[risultatoIO] (345.60,198.91) -- (393.23,198.91);
\node[anchor=north west, align=left] at (357.87,224.02)
  {Software\\validato};

% --- freccia grande della validazione ------------------------------------
% poligono a punta di freccia rivolto a sinistra: punta, testa, coda
\filldraw[fill=black!11, draw=black, line width=0.9pt]
  (164.08,197.25) -- (194.40,211.54) -- (194.40,205.40) --
  (257.05,205.40) -- (257.05,189.10) -- (194.40,189.10) --
  (194.40,182.96) -- cycle;
\node at (225.70,197.25) {Validazione};

% --- frecce dei risultati (continue) --------------------------------------
\draw[risultato] (92.37,166.19) -- (92.37,131.34) -- (110.76,131.34);
\draw[risultato] (121.69,114.87) -- (121.69,79.94) -- (140.00,79.94);
\draw[risultato] (153.59,63.54) -- (153.59,28.61) -- (171.98,28.61);
\draw[risultato] (246.94,28.69) -- (265.25,28.69) -- (265.25,63.54);
\draw[risultato] (278.84,80.01) -- (291.55,80.01) -- (291.55,114.87);
\draw[risultato] (307.34,131.04) -- (319.90,131.04) -- (319.90,166.12);

% --- frecce delle verifiche (punteggiate) ---------------------------------
% NB: le due orizzontali risalgono da destra a sinistra, dalla fase di test
%     alla corrispondente fase di progettazione
\draw[verifica] (121.76,147.73) -- (121.76,166.19);
\draw[verifica] (153.67,96.41) -- (153.67,114.87);
\draw[verifica] (183.35,45.01) -- (183.35,63.54);
\draw[verifica] (222.57,79.94) -- (196.28,79.94);
\draw[verifica] (243.17,131.34) -- (167.03,131.34);

% --- diagonali delle attivita' --------------------------------------------
\draw[diagonale] (52.26,153.27) -- (122.50,42.60);
\node[align=center] at (50.92,96.20) {Attività di\\progettazione};
\draw[diagonale] (283.42,42.28) -- (353.50,152.80);
\node[align=center] at (351.77,96.20) {Attività di\\test};

% --- legenda ---------------------------------------------------------------
\draw[cornice, fill=white] (13.35,16.06) rectangle (106.32,46.26);
\draw[risultato] (19.55,37.88) -- (37.42,37.88);
\node[anchor=west] at (46.87,37.15) {Risultato};
\draw[verifica] (19.55,23.22) -- (37.42,23.22);
\node[anchor=west] at (46.87,22.43) {Verifica};

% --- cornice esterna --------------------------------------------------------
\draw[cornice] (0.37,0.37) rectangle (413.73,292.79);

\end{tikzpicture}
